% Abhakseite „Das kann ich“ – Zone nur im Gesamt (\mitzone)
%% TODO Ich-kann-Sätze: Platzhalter wie in den Aufgabendateien
\begin{abhakseite}
\ifdefined\mitzone
\abhakgruppe{Kennst du schon}
\abhak{1}{Bruch als Anteil lesen und kürzen (drei von zwölf, gekürzt ein Viertel), auf den Nenner hundert erweitern}
\abhak{2}{Bruch ↔ Dezimalzahl (ein Viertel, drei Fünftel)}
\abhak{3}{Durch hundert teilen, mit Dezimalzahlen multiplizieren (Kommaverschiebung)}
\abhak{4}{Hoch- und Runterrechnen in einer Tabelle (Dreisatz, proportionale Zuordnung)}
\abhak{5}{Runden auf eine Dezimalstelle}
\abhak{6}{Bruchteil einer Größe (zwei Drittel von 60 €)}
\abhak{7}{Bruchteil einer Größe (zwei Drittel von 60 €) – Fehler finden}
\abhak{8}{Bruchteil einer Größe (zwei Drittel von 60 €) – selbst rechnen}
\fi
\abhakgruppe{Einheit 1 · Prozente als Anteile}
\abhak{9}{Umwandeln}
\abhak{10}{Prozentangaben ordnen}
\abhak{11}{Umwandeln – Fehler finden, Begründen, Darstellungswechsel, Anwendung}
\abhakgruppe{Einheit 2 · Prozentsatz berechnen}
\abhak{12}{Was ist das Ganze?}
\abhak{13}{Streifen einteilen}
\abhak{14}{Prozentsatz}
\abhak{15}{Umkehrung: zu einem Prozentsatz ein Zahlenpaar angeben}
\abhak{16}{Prozentsatz – Fehler finden, Begründen, Darstellungswechsel, Anwendung}
\abhakgruppe{Einheit 3 · Prozentwert berechnen}
\abhak{17}{Prozentwert}
\abhak{18}{Mehrwertsteuer in Euro}
\abhak{19}{Prozentwert – Fehler finden, Begründen, Darstellungswechsel, Anwendung}
\abhakgruppe{Einheit 4 · Grundwert berechnen}
\abhak{20}{Was ist gegeben, was gesucht?}
\abhak{21}{Grundwert}
\abhak{22}{Grundwert – Fehler finden, Begründen, Darstellungswechsel, Anwendung}
\abhakgruppe{Einheit 5 · Prozentuale Veränderung}
\abhak{23}{Welcher Wert ist der alte?}
\abhak{24}{Veränderung}
\abhak{25}{Veränderung – Fehler finden, Begründen, Anwendung}
\end{abhakseite}
